\section{状态、行动、损失与风险构成一个决策问题}
\label{sec:11-Mathematical-Statistics/07-Statistical-Decision-and-Reliable-Prediction/01}

上线评审会的那一页写着：AUC \(0.94\)，准确率 \(98.7\%\)，两项都比上一版好。会开了四十分钟，通过。三个月后复盘，模型没有退化，损失却比预算高出一截。翻账才发现：漏掉一笔欺诈交易，平均赔出去 \(8000\) 元；误拦一笔正常交易，代价是一通确认电话，算到 \(20\) 元。两者差 \(400\) 倍。而线上用的判定阈值是 \(0.5\)。

那个 \(0.5\) 对应的假设是两种错误一样贵。没有人在任何文档里写下过这条假设，它是默认值——不写损失函数时，系统替你填进去的那一个。这一节把决策问题的骨架摆开，好让这类默认值无处藏身。

\subsection{三件套：未知状态、可选行动、损失}

一个决策问题需要三样东西才立得住。第一是未知状态 \(\theta\in\Theta\)：这笔交易是否欺诈，这个参数的真值是多少，下周的需求量是几。第二是可选行动 \(a\in\mathcal A\)：放行、拦截、转人工；报一个点估计；给一个区间。第三是损失 \(L(\theta,a)\)，说明状态为 \(\theta\) 而你做了 \(a\) 时要付多少代价。再补上数据 \(X\sim P_\theta\)，以及把数据映成行动的规则 \(\delta:\mathcal X\to\mathcal A\)，问题就完整了。

行动由数据决定，数据随机，因此 \(L(\theta,\delta(X))\) 是随机变量。对数据取期望，得到频率学派风险

\[
R(\theta,\delta)=\E_\theta\!\left[L\bigl(\theta,\delta(X)\bigr)\right].
\]

下标 \(\theta\) 提醒一件容易滑过去的事：期望是在真值为 \(\theta\) 的那个分布下取的。所以 \(R\) 不是一个数，是关于 \(\theta\) 的一整条曲线。样本量、噪声水平、规则本身的保守程度，全被压进这条曲线的形状里。要把曲线收成一个数，还得再选一次平均方式：按先验 \(\pi\) 加权，得到 Bayes 风险 \(r(\pi,\delta)=\int R(\theta,\delta)\,\pi(\mathrm d\theta)\)；或者取上确界 \(\sup_\theta R(\theta,\delta)\)，得到最坏情形风险。这两条路通向不同的规则，本单元第三节专门比较。

\subsection{估计、检验与预测只是行动集合换了一个}

点估计里 \(\mathcal A=\Theta\)。平方损失 \((a-\theta)^2\) 的风险就是 MSE，绝对损失 \(\abs{a-\theta}\) 把最优行动挪到中位数。更实用的是 pinball 损失

\[
L(\theta,a)=\tau\,(\theta-a)_+ + (1-\tau)\,(a-\theta)_+,
\]

它的最优行动是 \(\tau\) 分位数。备货备到需求的 \(0.9\) 分位数，等价于宣布缺货比压货贵 \(9\) 倍——容量规划里那句``按 P90 备''，说的就是这个损失比。

假设检验里 \(\mathcal A=\{0,1\}\)，损失是一个两格的表。固定显著性水平相当于给风险曲线的一段加了硬上界，再在这个约束下尽量抬高功效，这是\hyperref[sec:11-Mathematical-Statistics/05-Hypothesis-Testing/01]{``两类错误与功效共同定义检验''}里那套安排的决策论读法。预测里，行动可以是点、概率、集合，也可以是``不作答，转人工''，后两种分别是本单元第五、第六节的主题。

换一个损失，最优行动就换一个。同一份后验分布能支撑均值、中位数、分位数、集合等多种输出，它们不是同一个答案的不同格式。所以``哪个估计量最好''在损失写下来之前没有答案。

\subsection{损失不写下来，等于默认选了 0-1 损失}

分类里那条几乎不用想的判据 \(\hat y=\argmax_y P(y\given x)\)，对应的损失是 \(L(y,a)=\ind\{a\ne y\}\)：错就是错，两格代价都是 \(1\)。二分类时它给出阈值 \(0.5\)。把开头那两个数字填进去，误报代价 \(c_{\mathrm{FP}}=20\)、漏检代价 \(c_{\mathrm{FN}}=8000\)，最小化期望损失的判据变成拦截当且仅当

\[
P(y=1\given x)\ \ge\ \frac{c_{\mathrm{FP}}}{c_{\mathrm{FP}}+c_{\mathrm{FN}}}=\frac{20}{8020}\approx 0.0025 .
\]

推导留到下一节，这里只看结论：合适的阈值与 \(0.5\) 差了两个数量级。评审会上那个模型不是概率估错了，是概率算完之后那一步没人管。

由此得到一条可以直接执行的检查项：把损失矩阵的几格数字写在报告首页。填不出来，说明这次上线要买的东西还没想清楚；填得出来，阈值、报警量、人工复核规模全都跟着确定。写不写这几个数，决定权其实不在你手上——系统总要按某个损失运行，区别只在于它是你选的还是它自己填的。

\subsection{一条规则被处处压住，就该丢掉}

规则之间怎么比？把它们的风险曲线画在同一张图上。

\begin{center}
\begin{tikzpicture}[font=\footnotesize,>=stealth,x=1.25cm,y=1.25cm]
  \draw[->,black!70] (-0.15,0) -- (5.4,0) node[below] {\(\theta\)};
  \draw[->,black!70] (0,-0.12) -- (0,3.6) node[above] {\(R(\theta,\delta)\)};
  \draw[black!85,thick] plot[domain=0:4.6,samples=2] (\x,{1.6});
  \draw[black!85,thick,dashed] plot[domain=0:4.6,samples=60] (\x,{0.4+0.1*\x*\x});
  \draw[black!45,thick,dotted] plot[domain=0:4.6,samples=60] (\x,{0.9+0.12*\x*\x});
  \draw[black!35] (3.46,0) -- (3.46,1.6);
  \node[black!55,anchor=north] at (3.46,-0.05) {\(\theta^\ast\)};
  \node[black!85,anchor=west] at (4.65,1.6) {\(\delta_1\)};
  \node[black!85,anchor=west] at (4.65,2.52) {\(\delta_2\)};
  \node[black!55,anchor=west] at (4.65,3.44) {\(\delta_3\)};
  \node[black!60,anchor=west] at (1.15,0.2) {\(\delta_2\) 在原点附近更省};
  \node[black!60,anchor=west] at (1.5,3.15) {\(\delta_3\) 处处高于 \(\delta_2\)};
\end{tikzpicture}

\emph{三条风险曲线。\(\delta_3\) 被 \(\delta_2\) 整条压住，可以直接淘汰；\(\delta_1\) 与 \(\delta_2\) 在 \(\theta^\ast\) 交叉，谁更好取决于 \(\theta\) 落在哪一侧。}
\end{center}

图里的虚线与点线之间没有悬念：对每个 \(\theta\) 都有 \(R(\theta,\delta_2)\le R(\theta,\delta_3)\)，且至少一处严格小于，这时说 \(\delta_2\) 支配 \(\delta_3\)，而 \(\delta_3\) 是不可容许的。在既定的模型与损失下，保留它没有理由。

实线与虚线之间就有悬念了。\(\delta_2\) 是收缩型规则，把估计往原点拉，参数确实小的时候占便宜，参数大的时候要付偏差的账；\(\delta_1\) 的风险不随 \(\theta\) 变。两条线交叉，意味着``哪个更好''这个问题在信息上不完整——还得说清 \(\theta\) 大概落在哪个范围，或者你更怕平均表现差，还是更怕最坏情形差。只报一句``平均误差更小''会把这些遮住：平均是按哪个状态分布算的，测试集又是怎样组成的。

\subsection{模型、损失与规则是三层，失败要先归层}

模型 \(P_\theta\) 说数据怎样随机，损失 \(L\) 说错误怎样计价，规则 \(\delta\) 说看到数据后怎样动。三者可以分别出错，修法完全不同。

住院风险模型把概率估得很准，阈值却停在 \(0.5\)，问题在决策层，改一行常数即可。标签系统把没有复诊的病人一律记成阴性，问题在数据与模型层，阈值调到哪里都救不回来。部署人群从体检中心换成急诊科，则是 \(\theta\) 的分布变了，历史风险曲线不再描述当前风险，这一类见\hyperref[sec:13-Machine-Learning/01-Learning-Problems-and-Evaluation/04]{``数据划分、泄漏与分布偏移''}。把三层混着讲成``模型不好用''，可以改三个月还改不对地方。

\subsection{损失是价值判断，数学不替你做}

损失函数能把价值取舍写清楚，却不能从数据里发现什么后果更重要。金钱、健康、公平、不可逆的伤害压成一个标量，这一步是承诺，不是推导。同一个系统在不同人眼里损失矩阵不同：对平台是赔付金额，对被误拦的用户是一次交易失败加一通电话，两者不共享单位。即便必须行动，报告里也该写清损失是谁的视角、代价能不能补偿、哪些约束没有进目标函数。

数学保证的是另一件事：状态、模型、行动与损失一旦明确给定，规则之间就能被一致地比较。它不保证这套给定本身描述了现实。

损失写下来之后，剩下的问题是怎么用它。下一节把这一步拆成两段——先算后验概率，再按损失比取阈值——并给出那个阈值的公式和一个具体算例。
